% D04-010. AI-04 P0291--P0330, сохранённый article_report.
\section{Выявление насыщения измерительных трансформаторов тока}
\label{sec:ch4_ct_saturation}
\conceptsource{Статья 4, раздел 7. Перенос общего физического
кодировщика на другую задачу РЗА, формирование метки по погрешности,
позонные и файловые результаты, реальные примеры и ограничение
бинарного описания непрерывного искажения. Числа сверены с
сохранённым отчётом опыта CT saturation v2.}

Рассмотренные методы могут использоваться не только
для определения вида события, но и для анализа
качества поступающих измерений. При насыщении ТТ
вторичный ток перестаёт точно воспроизводить первичный.
Это искажает рассчитанные величины и может изменять
решение измерительного органа.
Для БАВР такой анализ представляет интерес
как дополнительная информация о достоверности
измерительных цепей.

В этой задаче сохраняется общий принцип Physical KAN-Transformer,
но меняются входные каналы, продолжительность наблюдения
и целевой выход. Отдельное сравнение всех архитектур
здесь не повторяется. Проверяется применение
физически мотивированной обработки к другой
измерительной задаче.

\subsection{Данные и формирование целевых меток}
\label{sec:ch4_ct_targets}
\conceptsource{Статья 4, экспериментальная база насыщения ТТ.
Идеальный ток используется для целевой разметки и контрольной
разности, но не подаётся на вход модели.}

Использовано 115\,971 математически смоделированное
наблюдение. Записи содержат идеальные первичные токи,
вторичные токи измерительных трансформаторов
и напряжения трёх фаз. Для целевой разметки используются предоставленные
с расчётным набором флаги искажения по фазам.
Они выделяют участки расхождения вторичного тока
с эталонным сигналом. По этим меткам оценивается
распознавание искажённого измерения.
\conceptsource{Уточнение R4-13 остаётся у источника данных.
Автор предполагает порог 10 процентов от максимальной амплитуды
эталонного сигнала конкретного опыта, но формула коллегами
пока не подтверждена. Это предположение не использовано
для пересчёта меток. Повторять запрос до появления источника не требуется.}

% D04-011 R4-13 BEGIN
\iffalse
\par\smallskip\begingroup\small\raggedright\color{red!60!black}
\noindent\textbf{Примечание рецензента [R4-13], существенно.} Не определён знаменатель этих 10 процентов: мгновенный идеальный ток, его амплитуда, номинальный ток или иной масштаб. Это меняет смысл метки, особенно около перехода тока через нуль. Местный загрузчик читает готовые flag\_sat\_phsA/B/C, а не воспроизводит критерий. Уточнить определение в исходной модели или прямо назвать это предоставленной модельной разметкой с пока не восстановленным критерием; не изобретать формулу.\par\endgroup\smallskip
\fi
% D04-011 R4-13 END
\iffalse
\par\smallskip\begingroup\small\raggedright\color{blue!55!black}
\noindent\textbf{Ответ на R4-13, D04-012-B.} В основном тексте используются предоставленные флаги модельной разметки. Предположение автора о процентах от максимальной амплитуды сохранено служебно и не объявлено восстановленным критерием. Запрос источника остаётся отложенным.\par\endgroup\smallskip
\fi


На вход сети поступают только вторичные токи
и напряжения, доступные в измерительном тракте.
Идеальный первичный ток используется при подготовке
целевого ответа и для сопоставления с измеренным.
Он не включается в признаки.
Выход образуют три независимые оценки
для ТТ фаз A, B и C.

В обучающей части содержится 93\,063 записи,
в проверочной 22\,908. Имена файлов этих частей
не пересекаются. При этом близкие точки
параметрической сетки моделирования отдельно
по группам не разделялись.
Следовательно, проверка охватывает новые файлы,
но не гарантирует независимость всех близких
по физическим параметрам режимов.

\subsection{Представление сигналов и настройка модели}
\label{sec:ch4_ct_inputs}
\conceptsource{Статья 4 и сохранённая конфигурация опыта.
Окно два периода, полярный вход, отдельные три выхода.
Масштабирование здесь 100 А/300 В не отождествляется
с делителями других экспериментальных серий.}

Вторичные токи делятся на 100~А,
напряжения на 300~В. Эти делители соответствуют
номинальным значениям 5~А и 100~В
с коэффициентами запаса 20 и 3.
Так сохраняется общий способ нормирования
с учётом масштаба измерения, но сами делители
задаются для данной задачи.

Наблюдаемое сетью окно сокращено до двух периодов.
Для сети 50~Гц это 40~мс.
Спектральные векторы размещаются с шагом
в одну тридцать вторую периода,
поэтому входное окно содержит 64 временных положения.
Используются первые девять гармоник,
дополнительные низкочастотные компоненты
с окнами 2, 4, 6 и 10 периодов
и симметричные составляющие основной частоты.
В результате сформированы 168 входных координат.

Дополнительные спектральные компоненты рассчитываются
по предшествующим измерениям.
Поэтому окно сети длительностью 40~мс
не описывает всю необходимую предысторию
спектрального преобразования.
При начале регистрации дополнительно требуется
накопить достаточный материал для этих расчётов.

Скрытая размерность составляет 48,
число голов внимания 4, число повторяемых слоёв 6.
Для перестановки фаз синхронно меняются токи,
напряжения и соответствующие целевые метки.
Дополнительно применяется общее изменение масштаба
от 0,9 до 1,1, не меняющее взаимное соответствие каналов.

\subsection{Результаты на моделируемых данных}
\label{sec:ch4_ct_results}
\conceptsource{Статья 4, позонная/файловая проверка и рисунки 11--12.
Основные цифры относятся к исходной фазе A. Временное
совпадение склеенной кривой не является измеренным временем
действия потокового терминала.}

В таблице~\ref{tab:ch4_ct_metrics} приведены
результаты для фазы A при пороге 0,5. Основной отчёт
рассчитывает показатели по первому фазному выходу.
Они не усредняются по трём выходам и не характеризуют
их одинаково полную проверку. Перестановки фаз
рассматриваются как отдельный дополнительный опыт.
На уровне временных положений оценивается
совпадение каждого целевого признака.
На уровне файла насыщение считается обнаруженным,
если положительный выход появился хотя бы один раз.
Именно различие этих единиц объясняет
заметно разные показатели.

% D04-011 R4-14 BEGIN
\ifdefined\AuthorReview
\par\smallskip\begingroup\small\raggedright\color{red!60!black}
\noindent\textbf{Примечание рецензента [R4-14], уточнить.} Ограничение фазой A указано правильно, но читателю нужно объяснить его причину: основной независимый показатель не охватывает три выхода одинаково. Есть отдельный phase\_transfer.json с проверкой перестановок фаз; если он включается, обозначить его как отдельную проверку, а не расширять основную таблицу без проверки выборки. Сами числа Precision/Recall/F1 и файловая матрица здесь совпадают с сохранённым metrics.json.\par\endgroup\smallskip
\fi
% D04-011 R4-14 END
\ifdefined\AuthorReview
\par\smallskip\begingroup\small\raggedright\color{blue!55!black}
\noindent\textbf{Ответ на R4-14, D04-012-B.} Указан первый фазный выход основного отчёта и отсутствие усреднения трёх выходов. Отдельный phase\_transfer.json не подмешан в основную таблицу. Числа основной таблицы не менялись.\par\endgroup\smallskip
\fi

\iffalse
(Примечание автора: Пусть рецензент тут ответит, вроде норм.) **[AA119: рассмотрено 29.10; ответ ниже]**
\fi
\authorreviewnote{AA119}

\begin{table}[htbp]
 \centering
 \caption{Распознавание насыщения ТТ фазы A
 на проверочных моделируемых записях}
 \label{tab:ch4_ct_metrics}
 \begin{tabular}{lrrr}
  \hline
  Единица оценки & Precision & Recall & F1 \\
  \hline
  Временное положение & 0,667 & 0,695 & 0,681 \\
  Осциллограмма & 0,993 & 0,814 & 0,895 \\
  \hline
 \end{tabular}
\end{table}

Покомпонентные числа показаны на
рисунке~\ref{fig:ch4_ct_confusion}.
Отрицательные временные положения значительно
преобладают. Поэтому высокая общая доля
правильных ответов скрыла бы часть пропусков
и лишних положительных выходов.
Для рассматриваемой задачи основными
оставлены Precision, Recall и F1.

\begin{figure}[htbp]
 \centering
 \includegraphics[width=0.76\textwidth]{ch04_ct_confusion_article.png}
 \caption{Матрица ошибок позонного выявления
 насыщения ТТ фазы A}
 \label{fig:ch4_ct_confusion}
\end{figure}
\figurereview{Проверить перевод подписей матрицы и читаемость чисел. Числа соответствуют исходной оценке, изменять их при оформлении не требуется.}

На файловом уровне получено 5802 обнаружения
при 1328 пропущенных положительных записях.
В отрицательной группе зарегистрировано 38 ложных
обнаружений и 15\,740 правильных отрицательных ответов.
Для постоянной времени апериодической составляющей
0,015~с файловый F1 составляет 0,918,
для 0,4~с составляет 0,873.
Таким образом, наличие события в записи определяется
лучше, чем весь размеченный интервал искажения.

На рисунке~\ref{fig:ch4_ct_simulation}
показан отдельный моделируемый опыт.
Вместе с вторичными токами приведена разность
между идеальным и измеренным током,
приведёнными к одному масштабу.
Изменение выходной оценки сопоставляется
с величиной этого искажения.

\begin{figure}[p]
 \centering
 \includegraphics[width=\textwidth]{ch04_ct_simulation_article.png}
 \caption{Токи, контрольная разность и выходы
 модели в опыте с насыщением ТТ}
 \label{fig:ch4_ct_simulation}
\end{figure}
\figurereview{Авторская проверка рисунка ещё требуется. Сохранить различие токов, контрольной разности и выходной оценки, проверить подписи и единицы.}

Для отображения одной кривой по записи
оценки перекрывающихся окон объединяются.
Для оценки временного совпадения ретроспективной разметки
медианное расстояние
между первой истинной меткой и первым последующим
обнаружением составляет 1,24~мс,
95-й процентиль составляет 20,46~мс.
Оценка относится к 5802 обнаруженным
положительным записям и не включает пропуски
как конечные задержки.

Для каждой записи выбирается первая истинная метка
и первый положительный ответ не раньше неё.
Если ответ возник заранее и сохраняется на этой границе,
получается нулевой сдвиг. Ранние положительные участки
сами по себе не штрафуются этой величиной.
Она не характеризует отсутствие ложного пуска
или устойчивость нового состояния.

% D04-011 R4-15 BEGIN
\ifdefined\AuthorReview
\par\smallskip\begingroup\small\raggedright\color{red!60!black}
\noindent\textbf{Примечание рецензента [R4-15], уточнить.} Оговорка о ретроспективности корректна и должна остаться рядом с числом. Дополнительно уточнить обработку уже положительного выхода до начала целевой метки и критерий первого последующего обнаружения. Медиана по найденным записям не учитывает 1328 пропусков и не характеризует устойчивость выхода. Для окончательного текста полезно назвать величину «временное совпадение ретроспективной разметки», сохранив её отдельно от быстродействия защиты.\par\endgroup\smallskip
\fi
% D04-011 R4-15 END
\ifdefined\AuthorReview
\par\smallskip\begingroup\small\raggedright\color{blue!55!black}
\noindent\textbf{Ответ на R4-15, D04-012-B.} Добавлены критерий первого ответа не раньше истинной метки и нулевой сдвиг при продолжающемся раннем ответе. Величина названа временным совпадением ретроспективной разметки; 1328 пропусков не скрыты, оперативной задержкой она не называется.\par\endgroup\smallskip
\fi

(Примечание автора: Пусть рецензент тут ответит, вроде норм.)

Поскольку объединение использует оценки
разных окон для одного временного положения,
это характеристика временного совпадения
при анализе зарегистрированного материала.
Она не заменяет проверку задержки
оперативного алгоритма, использующего
только уже накопленные к моменту решения данные.

\subsection{Примеры реальных записей и границы интерпретации}
\label{sec:ch4_ct_real_examples}
\conceptsource{Статья 4, реальные примеры рисунка 13 и manifest
исходных рисунков. Реальный архив не имеет эталонных меток
насыщения, поэтому изображения рассматриваются качественно.}

Дополнительное применение выполнено на 981 группе
реальных осциллограмм, ранее подготовленных
для исследований переходных процессов.
В этом материале нет эталонного идеального тока
и независимой разметки насыщения ТТ.
Поэтому по нему не рассчитывается точность
обнаружения физического насыщения.

В примере на рисунке~\ref{fig:ch4_ct_real_switching}
наблюдается бросок тока с выраженной
апериодической составляющей.
На основной части процесса выходы модели
не превышают порог 0,5.
Это показывает её реакцию на данный пример,
но само по себе не устанавливает отсутствие
физического насыщения измерительного трансформатора.

\begin{figure}[htbp]
 \centering
 \includegraphics[width=\textwidth]{ch04_ct_real_switching_article.png}
 \caption{Выходы модели на реальном
 коммутационном процессе с броском тока}
 \label{fig:ch4_ct_real_switching}
\end{figure}
\figurereview{Авторская проверка рисунка ещё требуется. Перевести легенду и сохранить отметку реального коммутационного процесса.}

На другом фрагменте,
показанном на рисунке~\ref{fig:ch4_ct_real_candidate},
оценка фазы A около 240~мс достигает 0,728.
Здесь заметно локальное изменение формы
измеренного тока.
Такой интервал выделяется как кандидат
для дальнейшего предметного рассмотрения.
Положительный ответ сети не используется
вместо недостающего физического эталона.

\begin{figure}[htbp]
 \centering
 \includegraphics[width=\textwidth]{ch04_ct_real_candidate_article.png}
 \caption{Реальный фрагмент, выделенный моделью
 для последующей проверки насыщения ТТ}
 \label{fig:ch4_ct_real_candidate}
\end{figure}
\figurereview{Авторская проверка рисунка ещё требуется. Укрупнить подписи; выход модели на реальной записи не заменяет её экспертную разметку.}

\subsection{Непрерывное искажение и двоичная метка}
\label{sec:ch4_ct_continuous_error}
\conceptsource{Статья 4, обсуждение границы10\% и дальнейшей
постановки. Различаются проведённая классификация и возможные
задачи восстановления тока/оценки погрешности, для которых
новые результаты здесь не заявляются.}

Физическая погрешность изменяется непрерывно,
а целевой признак переключается при достижении
принятой границы 10\%.
Близкие по величине отклонения могут оказаться
по разные стороны этой границы.
Это усложняет интерпретацию небольших ошибок
в начале и конце размеченного интервала.

В проведённом опыте сохраняется заданная
двоичная постановка. Дальнейшее развитие
может включать оценку самой погрешности
или восстановление неискажённого тока
с последующим применением инженерного порога.
Тогда непрерывный выход был бы непосредственно
связан с измерительной величиной.
Эти задачи не считаются уже решёнными
полученными результатами классификации.

Таким образом, общий способ обработки применён
к двум разным назначениям.
При ОЗЗ он распознаёт признаки переходного процесса,
при насыщении ТТ выделяет возможное искажение измерения.
Для дальнейшего использования в БАВР
эти выходы рассматриваются совместно
с назначением органа и качеством его входных сигналов.
Ни один из них отдельно не задаёт полную
логику переключения питания.


% D04-011 R4-17 BEGIN
\ifdefined\AuthorReview
\par\smallskip\begingroup\small\raggedright\color{red!60!black}
\noindent\textbf{Примечание рецензента [R4-17], важно.} Глава заканчивается локальным итогом опыта ТТ, но общего вывода по главе пока нет. Предлагаю после согласования кратко сформулировать: вклад автора в обработку и физические блоки; подтверждённый результат выбора входов; результаты и ограничения детекции на реальных данных; что передаётся в методику главы 5. Не повторять описание всех сетей. Отдельно сохранить честный результат: в текущей серии физических моделей нет единственного победителя по чувствительности и ложным обнаружениям.\par\endgroup\smallskip
\fi
% D04-011 R4-17 END
\ifdefined\AuthorReview
\par\smallskip\begingroup\small\raggedright\color{blue!55!black}
\noindent\textbf{Ответ на R4-17, D04-012-B.} Добавлен общий раздел выводов главы. Акцент согласован с автором — развитие архитектур, влияние представления и временного контекста, проверка на реальных задачах и переход к измерительному органу. Новый текст ожидает вычитки.\par\endgroup\smallskip
\fi
